%% toc: P3 32 32.5,32.6,32.7
\section{Proofs for the stability spectrum and the separating theories}
\label{sec:annex-p03-w03-p03-032-spectrum}

Throughout, \(T\) is complete and all tuples and sets lie in the ambient
model \(\mathfrak C\) of \cref{def:p03-w03-p03-031-monster}, whose saturation
cardinal exceeds \(2^\lambda\) for every cardinal \(\lambda\) under
discussion.  By \cref{thm:p03-w03-p03-031-saturation}\,(1), a set of
formulas in a small number of variables with parameters from a small set is
realized in \(\mathfrak C\) as soon as each of its finite subsets is; this is
called \emph{realization by saturation} below.

\begin{annexitem}{annex:p03-w03-p03-032-rank-basics}{thm:p03-w03-p03-032-rank-basics}{Monotonicity, finite character and splitting}
\begin{proof}
(1) If \(\rho\) is inconsistent, \(\mathrm R_\varphi(\rho)=-1\) and there is
nothing to prove.  Otherwise a \(\varphi\)-tree of height \(n\) above
\(\rho\) is one above \(\pi\): for each branch \(\eta\), a realization of
\(\rho\cup\{\varphi^{\eta(k)}(x;b_{\eta\mathbin{\upharpoonright}k}):k<n\}\)
realizes the smaller set with \(\pi\) in place of \(\rho\).  So every height
attained above \(\rho\) is attained above \(\pi\).

(2) One direction is (1).  Conversely assume that every finite
\(\pi_0\subseteq\pi\) has a \(\varphi\)-tree of height \(n\).  Take new
variables \(y_s\) for \(s\in2^{<n}\), of the length of \(y\), and \(x_\eta\)
for \(\eta\in2^n\), of the length of \(x\), and let \(\Sigma\) be the set of
formulas
\[
  \pi(x_\eta),\qquad
  \varphi^{\eta(k)}(x_\eta;y_{\eta\mathbin{\upharpoonright}k})
  \qquad(\eta\in2^n,\ k<n).
\]
A finite subset of \(\Sigma\) mentions only a finite \(\pi_0\subseteq\pi\);
a tree \((b_s)\) of height \(n\) above \(\pi_0\), together with realizations
\(c_\eta\) of its branches, satisfies that subset with \(y_s\mapsto b_s\),
\(x_\eta\mapsto c_\eta\).  There are finitely many variables, so \(\Sigma\)
is realized by saturation, and the values of the \(y_s\) form a
\(\varphi\)-tree of height \(n\) above \(\pi\).

For the second claim let \(k=\mathrm R_\varphi(\pi)<\infty\).  There is no
tree of height \(k+1\) above \(\pi\) (for \(k=-1\): \(\pi\) is inconsistent,
which is the absence of a tree of height \(0\)).  By the first claim some
finite \(\pi_0\subseteq\pi\) has no tree of height \(k+1\), so
\(\mathrm R_\varphi(\pi_0)\leq k\); by (1),
\(\mathrm R_\varphi(\pi_0)\geq k\).  If \(k\geq0\), then
\(\mathrm R_\varphi(\pi_0)\neq-1\) because \(\pi_0\subseteq\pi\) is
consistent; if \(k=-1\), the finite \(\pi_0\) is inconsistent.  In either
case \(\mathrm R_\varphi(\pi_0)=k\).

(3) For \(n=0\) the hypothesis says that \(\pi\cup\{\varphi(x;b)\}\) is
consistent, and the one-node family \(b_\varnothing=b\) is a tree of height
\(1\) above \(\pi\) because both \(\pi\cup\{\varphi(x;b)\}\) and
\(\pi\cup\{\neg\varphi(x;b)\}\) are consistent.  For \(n\geq1\) let
\((b'_s:s\in2^{<n})\) be a tree above \(\pi\cup\{\neg\varphi(x;b)\}\) and
\((b''_s:s\in2^{<n})\) one above \(\pi\cup\{\varphi(x;b)\}\).  Define, for
\(s\in2^{<n+1}\),
\[
  b_\varnothing=b,\qquad b_{\langle0\rangle^\frown s}=b'_s,\qquad
  b_{\langle1\rangle^\frown s}=b''_s .
\]
A branch \(\eta\in2^{n+1}\) has the form \(\langle\varepsilon\rangle^\frown
\eta'\) with \(\eta'\in2^n\).  For \(\varepsilon=0\) the set
\(\pi\cup\{\varphi^{\eta(k)}(x;b_{\eta\mathbin{\upharpoonright}k}):k\leq n\}\)
equals \(\pi\cup\{\neg\varphi(x;b)\}\cup
\{\varphi^{\eta'(k)}(x;b'_{\eta'\mathbin{\upharpoonright}k}):k<n\}\), which is
consistent because \((b'_s)\) is a tree; the case \(\varepsilon=1\) uses
\((b''_s)\).  So \((b_s)\) is a tree of height \(n+1\) above \(\pi\).
\end{proof}
\end{annexitem}

\begin{annexitem}{annex:p03-w03-p03-032-local-character}{thm:p03-w03-p03-032-local-character}{A finite part decides a \(\varphi\)-type}
\begin{proof}
Since \(p\) is consistent and \(\varnothing\subseteq p\),
\(0\leq\mathrm R_\varphi(p)\leq\mathrm R_\varphi(x=x)<\infty\) by
\cref{thm:p03-w03-p03-032-rank-basics}\,(1), so \(k\) is a natural number
and a finite \(p_0\subseteq p\) with \(\mathrm R_\varphi(p_0)=k\) exists by
part~(2) of that lemma.

Let \(b\) be a tuple from \(A\).  If \(\varphi(x;b)\in p\), then
\(p_0\cup\{\varphi(x;b)\}\subseteq p\), so its rank is at least
\(\mathrm R_\varphi(p)=k\) by monotonicity.  If \(\varphi(x;b)\notin p\),
then \(\neg\varphi(x;b)\in p\) by completeness, and in the same way
\(\mathrm R_\varphi(p_0\cup\{\neg\varphi(x;b)\})\geq k\).  Were also
\(\mathrm R_\varphi(p_0\cup\{\varphi(x;b)\})\geq k\), the splitting clause
\cref{thm:p03-w03-p03-032-rank-basics}\,(3) would give
\(\mathrm R_\varphi(p_0)\geq k+1\), contrary to the choice of \(p_0\).  This
proves the displayed equivalence.

The right-hand side of the equivalence depends only on \(p_0\) and \(b\),
because \(k=\mathrm R_\varphi(p_0)\).  Hence two complete \(\varphi\)-types
over \(A\) for which the same finite set \(p_0\) was chosen contain the same
instances \(\varphi(x;b)\), and therefore also the same instances
\(\neg\varphi(x;b)\); they are equal.  The chosen \(p_0\) thus gives an
injection of \(S_\varphi(A)\) into the set of finite sets of
\(\varphi\)-instances over \(A\).  There are at most
\(2\cdot\abs A^{\abs y}+\aleph_0\leq\abs A+\aleph_0\) instances, hence at most
\(\abs A+\aleph_0\) finite sets of them.
\end{proof}
\end{annexitem}

\begin{annexitem}{annex:p03-w03-p03-032-product}{thm:p03-w03-p03-032-product}{Types are determined by their local parts}
\begin{proof}
Every formula of a complete type \(p\in S_1(A)\) has the form
\(\psi(x;b)\) with \(\psi(x;y)\) a formula of \(\Lang\) and \(b\) from
\(A\), so it is an instance of \(\psi\).  Hence \(p\) is the union of the
sets \(p\mathbin{\upharpoonright}\varphi\), and
\(p\mapsto(p\mathbin{\upharpoonright}\varphi)_\varphi\) is an injection of
\(S_1(A)\) into \(\prod_\varphi S_\varphi(A)\).  Formulas are finite strings
over \(\Lang\) together with countably many variables, so there are at most
\(\abs T\) of them.
\end{proof}
\end{annexitem}

\begin{annexitem}{annex:p03-w03-p03-032-one-to-many}{thm:p03-w03-p03-032-one-to-many}{One variable suffices}
\begin{proof}
Let \(\abs A\leq\lambda\).  Define \(A_0=A\) and, given \(A_j\) with
\(\abs{A_j}\leq\lambda\), choose one realization in \(\mathfrak C\) of each
type in \(S_1(A_j)\) and let \(A_{j+1}\) be \(A_j\) together with these
realizations.  By hypothesis \(\abs{A_{j+1}}\leq\lambda\).

Let \(q(x_0,\ldots,x_{m-1})\in S_m(A)\).  We choose \(d_j\in A_{j+1}\) for
\(j<m\) such that, for every \(j\leq m\),
\begin{equation*}
  (\ast)_j\qquad \mathfrak C\entails
  \exists x_j\cdots\exists x_{m-1}\,
  \psi(d_0,\ldots,d_{j-1},x_j,\ldots,x_{m-1})
  \quad\text{for every }\psi\in q .
\end{equation*}
For \(j=0\) this holds because \(q\) is realized in \(\mathfrak C\).  Given
\((\ast)_j\) with \(j<m\), let \(\Pi\) be the set of formulas
\(\exists x_{j+1}\cdots\exists x_{m-1}\,\psi(d_0,\ldots,d_{j-1},x_j,\ldots)\)
for \(\psi\in q\), in the single variable \(x_j\) with parameters from
\(A_j\).  Since \(q\) is closed under conjunction up to equivalence, a finite
subset of \(\Pi\) is implied by one member of \(\Pi\), which is satisfiable
by \((\ast)_j\).  So \(\Pi\) is consistent and extends to some
\(r\in S_1(A_j)\), which is realized by some \(d_j\in A_{j+1}\); this
\(d_j\) satisfies \((\ast)_{j+1}\).  At \(j=m\), \((\ast)_m\) says that
\((d_0,\ldots,d_{m-1})\) realizes \(q\).  Distinct types have disjoint sets
of realizations, so \(q\mapsto(d_0,\ldots,d_{m-1})\) is injective, and
\(\abs{S_m(A)}\leq\abs{A_m}^m\leq\lambda\).
\end{proof}
\end{annexitem}

\begin{annexitem}{annex:p03-w03-p03-032-tree-count}{thm:p03-w03-p03-032-tree-count}{Trees of every height give many types}
\begin{proof}
\emph{Cardinals.}  Since \(2^\lambda>\lambda\), \(\mu\leq\lambda\).  Let
\(2^{<\mu}\) be the set of functions from ordinals \(\nu<\mu\) into
\(\{0,1\}\).  For \(\nu<\mu\) we have \(\abs\nu<\mu\), so
\(2^{\abs\nu}\leq\lambda\) by the minimality of \(\mu\), and
\(\abs{2^{<\mu}}\leq\mu\cdot\lambda=\lambda\).

\emph{A tree of height \(\mu\).}  Take variables \(y_s\) for
\(s\in2^{<\mu}\) and \(x_\eta\) for \(\eta\in2^\mu\), and let \(\Sigma\) be
the set of formulas
\(\varphi^{\eta(i)}(x_\eta;y_{\eta\mathbin{\upharpoonright}i})\) for
\(\eta\in2^\mu\) and \(i<\mu\).  Let \(\Sigma_0\subseteq\Sigma\) be finite;
it mentions finitely many \(\eta\) and finitely many levels, listed as
\(i_1<\cdots<i_r\).  By hypothesis there is a \(\varphi\)-tree
\((b_t:t\in2^{<r})\) of height \(r\) above the empty set, with a realization
\(c_\tau\) of each branch \(\tau\in2^r\).  Interpret
\[
  y_{\eta\mathbin{\upharpoonright}i_j}\mapsto
  b_{(\eta(i_1),\ldots,\eta(i_{j-1}))},
  \qquad
  x_\eta\mapsto c_{(\eta(i_1),\ldots,\eta(i_r))}.
\]
The first assignment is well defined: \((\eta(i_1),\ldots,\eta(i_{j-1}))\)
depends only on \(\eta\mathbin{\upharpoonright}i_j\) because
\(i_{j-1}<i_j\).  For \(\tau=(\eta(i_1),\ldots,\eta(i_r))\) the formula
\(\varphi^{\eta(i_j)}(x_\eta;y_{\eta\mathbin{\upharpoonright}i_j})\) becomes
\(\varphi^{\tau(j-1)}(c_\tau;b_{\tau\mathbin{\upharpoonright}(j-1)})\), which
holds because \(c_\tau\) realizes the branch \(\tau\).  So \(\Sigma\) is
finitely satisfiable; it has at most \(2^\mu\leq2^\lambda\) variables and no
parameters, so it is realized by saturation, by \(b_s\) and \(c_\eta\).

\emph{Counting.}  Let \(A\) be the set of entries of the \(b_s\); then
\(\abs A\leq\abs{2^{<\mu}}\cdot\abs y\leq\lambda\).  Let \(p_\eta\) be the
complete \(\varphi\)-type of \(c_\eta\) over \(A\).  If \(\eta\neq\eta'\) and
\(i\) is least with \(\eta(i)\neq\eta'(i)\), then
\(s=\eta\mathbin{\upharpoonright}i=\eta'\mathbin{\upharpoonright}i\),
\(\varphi^{\eta(i)}(x;b_s)\in p_\eta\) and
\(\varphi^{\eta'(i)}(x;b_s)\in p_{\eta'}\), and these are a formula and its
negation.  So \(p_\eta\neq p_{\eta'}\) and
\(\abs{S_\varphi(A)}\geq2^\mu>\lambda\).
\end{proof}
\end{annexitem}

\begin{annexitem}{annex:p03-w03-p03-032-erdos-makkai}{thm:p03-w03-p03-032-erdos-makkai}{Erd\H{o}s--Makkai theorem}
\begin{proof}
The proof follows Erd\H{o}s and Makkai \cite[Theorem, cases (i) and
(ii)]{ErdosMakkai1966}.  For \(B\subseteq I\) and a family
\(\mathcal H\) of subsets of \(I\), let
\(B\sqcap\mathcal H\defeq\{B\cap X:X\in\mathcal H\}\), the family of traces
of \(\mathcal H\) on \(B\); for \(F\subseteq I\), let
\(\mathcal H_F\defeq\{X\in\mathcal H:F\subseteq X\}\).  Complements are
taken in \(I\) and written \(\overline X=I\setminus X\).  Consider the
condition
\begin{quote}
\((\star)\) for every \(B\subseteq I\) and \(\mathcal H\subseteq\mathcal G\)
with \(\abs{B\sqcap\mathcal H}>m\), there is \(X\in\mathcal H\) with
\(\abs{(B\cap\overline X)\sqcap\mathcal H}>m\).
\end{quote}

\emph{Case 1: \((\star)\) holds.  Then \(\mathcal G\) strongly cuts a
sequence.}  We choose \(a_n\in I\) and \(X_n\in\mathcal G\) such that, for
every \(n\),
\begin{align*}
  &(\mathrm a)_n\quad a_0,\ldots,a_{k-1}\in X_k\text{ and }
     a_k,\ldots,a_{n-1}\notin X_k\text{ for all }k<n;\\
  &(\mathrm b)_n\quad
     \bigl\lvert(\overline{X_0}\cap\cdots\cap\overline{X_{n-1}})\sqcap
     \mathcal G_{\{a_0,\ldots,a_{n-1}\}}\bigr\rvert>m .
\end{align*}
For \(n=0\), \((\mathrm a)_0\) is empty and \((\mathrm b)_0\) says
\(\abs{I\sqcap\mathcal G}>m\), which holds because distinct subsets of \(I\)
have distinct traces on \(I\).  Assume \((\mathrm a)_n\) and
\((\mathrm b)_n\), and put
\(B=\overline{X_0}\cap\cdots\cap\overline{X_{n-1}}\) and
\(\mathcal H=\mathcal G_{\{a_0,\ldots,a_{n-1}\}}\).  By \((\star)\) there is
\(X_n\in\mathcal H\) with
\(\abs{(B\cap\overline{X_n})\sqcap\mathcal H}>m\).  Put
\(B'=B\cap\overline{X_n}\).  Some \(a_n\in B'\) satisfies
\(\abs{B'\sqcap\mathcal H_{\{a_n\}}}>m\): otherwise, since each nonempty
trace \(B'\cap X\) contains some \(a\in B'\) and then
\(X\in\mathcal H_{\{a\}}\),
\[
  B'\sqcap\mathcal H\subseteq
  \bigcup_{a\in B'}\bigl(B'\sqcap\mathcal H_{\{a\}}\bigr)\cup\{\varnothing\}
\]
would have at most \(m\cdot m+1=m\) members.  As
\(\mathcal H_{\{a_n\}}=\mathcal G_{\{a_0,\ldots,a_n\}}\), this is
\((\mathrm b)_{n+1}\).  For \((\mathrm a)_{n+1}\): \(a_n\in B'\) gives
\(a_n\notin X_k\) for all \(k\leq n\), and \(X_n\in\mathcal H\) gives
\(a_0,\ldots,a_{n-1}\in X_n\); the remaining clauses are \((\mathrm a)_n\).
Since every pair \(i,k\) occurs in some \((\mathrm a)_n\),
\(a_i\in X_k\Leftrightarrow i<k\) for all \(i,k\).

\emph{Case 2: \((\star)\) fails.  Then \(I\ominus\mathcal G\) strongly cuts a
sequence.}  Fix \(B\subseteq I\) and \(\mathcal H\subseteq\mathcal G\) with
\(\abs{B\sqcap\mathcal H}>m\) and
\(\abs{(B\cap\overline X)\sqcap\mathcal H}\leq m\) for every
\(X\in\mathcal H\).  Put \(I'=B\) and \(\mathcal G'=B\sqcap\mathcal H\), a
family of subsets of \(I'\) with \(\abs{I'}\leq m<\abs{\mathcal G'}\).
Complements in \(I'\) are written \(I'\setminus Y\).  For
\(Y=B\cap X\in\mathcal G'\) we have
\((I'\setminus Y)\sqcap\mathcal G'=(B\cap\overline X)\sqcap\mathcal H\), so
\begin{equation*}
  (\dagger)\qquad \abs{(I'\setminus Y)\sqcap\mathcal G'}\leq m
  \quad\text{for every }Y\in\mathcal G' .
\end{equation*}
We show that \((I',\ I'\ominus\mathcal G')\), where
\(I'\ominus\mathcal G'=\{I'\setminus Y:Y\in\mathcal G'\}\), satisfies
\((\star)\).  Let \(B'\subseteq I'\) and
\(\mathcal H'=\{I'\setminus Y:Y\in\mathcal K\}\) with
\(\mathcal K\subseteq\mathcal G'\) and \(\abs{B'\sqcap\mathcal H'}>m\).
Since \(B'\cap(I'\setminus Y)=B'\setminus(B'\cap Y)\), complementation in
\(B'\) is a bijection between \(B'\sqcap\mathcal H'\) and
\(B'\sqcap\mathcal K\), so \(\abs{B'\sqcap\mathcal K}>m\).  Pick any
\(Y_0\in\mathcal K\) and put \(X'=I'\setminus Y_0\in\mathcal H'\).  The
traces of \(\mathcal K\) on \(B'\setminus Y_0\) are images of its traces on
\(I'\setminus Y_0\) under \(Z\mapsto Z\cap B'\), so by \((\dagger)\) there are
at most \(m\) of them.  Every trace \(B'\cap Y\) is the union of
\(B'\cap Y_0\cap Y\) and \((B'\setminus Y_0)\cap Y\); if there were at most
\(m\) traces of \(\mathcal K\) on \(B'\cap Y_0\), there would be at most
\(m\cdot m=m\) traces on \(B'\), which is false.  So
\(\abs{(B'\cap Y_0)\sqcap\mathcal K}>m\).  Now \(B'\setminus X'=B'\cap Y_0\),
and complementation in \(B'\cap Y_0\) is a bijection between
\((B'\cap Y_0)\sqcap\mathcal K\) and \((B'\cap Y_0)\sqcap\mathcal H'\).  Hence
\(\abs{(B'\setminus X')\sqcap\mathcal H'}>m\), which is \((\star)\) for
\((I',I'\ominus\mathcal G')\).

Since \(\abs{I'\ominus\mathcal G'}=\abs{\mathcal G'}>m\geq\abs{I'}\), Case~1
applied to \((I',I'\ominus\mathcal G')\) gives \(b_i\in B\) and
\(Y_k=B\cap X_k\) with \(X_k\in\mathcal H\subseteq\mathcal G\) such that
\(b_i\in B\setminus Y_k\Leftrightarrow i<k\).  As \(b_i\in B\),
\(b_i\in B\setminus Y_k\) iff \(b_i\in I\setminus X_k\).  So the sets
\(I\setminus X_k\in I\ominus\mathcal G\) strongly cut \((b_i)_{i<\omega}\).
\end{proof}
\end{annexitem}

\begin{annexitem}{annex:p03-w03-p03-032-many-to-order}{thm:p03-w03-p03-032-many-to-order}{Many \(\varphi\)-types give the order property}
\begin{proof}
Let \(I\) be the set of tuples from \(A\) of the length of \(y\) and
\(m=\abs A+\aleph_0\); then \(\abs I\leq m\).  For \(p\in S_\varphi(A)\) put
\(w(p)=\{b\in I:\varphi(x;b)\in p\}\).  A complete \(\varphi\)-type is
determined by \(w(p)\), so \(\mathcal G=\{w(p):p\in S_\varphi(A)\}\) has more
than \(m\) members.  By \cref{thm:p03-w03-p03-032-erdos-makkai} there are
\(b_i\in I\) and types \(p_k\) such that either (i)
\(b_i\in w(p_k)\Leftrightarrow i<k\) for all \(i,k\), or (ii)
\(b_i\in I\setminus w(p_k)\Leftrightarrow i<k\), that is,
\(b_i\in w(p_k)\Leftrightarrow k\leq i\).  Let \(c_k\) realize \(p_k\).
Then \(\varphi(c_k;b_i)\) holds iff \(i<k\) in case (i), and iff \(k\leq i\) in
case (ii).

In case (ii), put \(c'_i=c_{i+1}\) and \(b'_j=b_j\): then
\(\varphi(c'_i;b'_j)\) holds iff \(i+1\leq j\) iff \(i<j\).  In case (i), let
\(\Sigma\) be the set of formulas \(\varphi(x_i;y_j)\) for \(i<j<\omega\)
and \(\neg\varphi(x_i;y_j)\) for \(j\leq i<\omega\).  A finite subset of
\(\Sigma\) mentions indices below some \(N\); it is satisfied by
\(x_i\mapsto c_{N-i}\) and \(y_j\mapsto b_{N-j}\), because
\(\varphi(c_{N-i};b_{N-j})\) holds iff \(N-j<N-i\) iff \(i<j\).  The set
\(\Sigma\) has countably many variables and no parameters, so it is realized
by saturation, by tuples \(c'_i,b'_j\) with \(\varphi(c'_i;b'_j)\)
iff \(i<j\).  In both cases \(\varphi\) has the order property.
\end{proof}
\end{annexitem}

\begin{annexitem}{annex:p03-w03-p03-032-order-to-tree}{thm:p03-w03-p03-032-order-to-tree}{The order property gives trees of every height}
\begin{proof}
Let \((c_i)_{i<\omega}\), \((b_j)_{j<\omega}\) satisfy
\(\varphi(c_i;b_j)\Leftrightarrow i<j\), and fix \(n<\omega\).  Assign to
each binary sequence \(s\) of length at most \(n\) a rational \(q_s\) by
\[
  q_\varnothing=\tfrac12,\qquad
  q_{s^\frown\langle\varepsilon\rangle}=q_s+(1-2\varepsilon)\,2^{-(\abs s+2)}
  \quad(\varepsilon\in\{0,1\}).
\]
By induction, \(q_s\) is an odd integer divided by \(2^{\abs s+1}\), so
\(q_s\neq q_t\) whenever \(\abs s\neq\abs t\).

\emph{Claim.} For \(\eta\in2^n\) and \(k<n\),
\(q_\eta<q_{\eta\mathbin{\upharpoonright}k}\) iff \(\eta(k)=1\).  Indeed
\(q_\eta-q_{\eta\mathbin{\upharpoonright}k}
=\sum_{j=k}^{n-1}(1-2\eta(j))2^{-(j+2)}\).  The term \(j=k\) has absolute
value \(2^{-(k+2)}\), and the other terms have absolute values summing to
\(\sum_{j=k+1}^{n-1}2^{-(j+2)}=2^{-(k+2)}(1-2^{-(n-k-1)})<2^{-(k+2)}\).  So
the difference has the sign of \(1-2\eta(k)\), which is negative exactly
when \(\eta(k)=1\).

Let \(d_0<d_1<\cdots<d_r\) list the finitely many rationals \(q_s\) with
\(\abs s\leq n\), and let \(f(d_t)=t\).  For \(s\in2^{<n}\) put
\(b'_s=b_{f(q_s)}\), and for \(\eta\in2^n\) put
\(c'_\eta=c_{f(q_\eta)}\).  Since \(\abs\eta=n\neq\abs s\), the values
\(q_\eta\) and \(q_s\) differ, so \(f(q_\eta)<f(q_s)\) iff \(q_\eta<q_s\).
For \(k<n\) and \(s=\eta\mathbin{\upharpoonright}k\), the Claim gives
\[
  \varphi(c'_\eta;b'_s)
  \iff f(q_\eta)<f(q_s)
  \iff q_\eta<q_{\eta\mathbin{\upharpoonright}k}
  \iff \eta(k)=1 .
\]
So \(c'_\eta\) realizes
\(\{\varphi^{\eta(k)}(x;b'_{\eta\mathbin{\upharpoonright}k}):k<n\}\), and
\((b'_s:s\in2^{<n})\) is a \(\varphi\)-tree of height \(n\) above the empty
set.  As \(n\) was arbitrary, \(\mathrm R_\varphi(x=x)=\infty\).
\end{proof}
\end{annexitem}

\begin{annexitem}{annex:p03-w03-p03-032-formula}{thm:p03-w03-p03-032-formula}{Equivalent forms of instability for one formula}
\begin{proof}
(1)\(\Rightarrow\)(2) is \cref{thm:p03-w03-p03-032-order-to-tree}, and
(2)\(\Rightarrow\)(3) is \cref{thm:p03-w03-p03-032-tree-count}.
(3)\(\Rightarrow\)(4): apply (3) with \(\lambda=\aleph_0\) to get \(A\) with
\(\abs A\leq\aleph_0\) and \(\abs{S_\varphi(A)}>\aleph_0=\abs A+\aleph_0\).
(4)\(\Rightarrow\)(1) is \cref{thm:p03-w03-p03-032-many-to-order}.  If the
four conditions fail, then (2) fails, so
\(\mathrm R_\varphi(x=x)<\infty\) and
\cref{thm:p03-w03-p03-032-local-character} gives
\(\abs{S_\varphi(A)}\leq\abs A+\aleph_0\).
\end{proof}
\end{annexitem}

\begin{annexitem}{annex:p03-w03-p03-032-spectrum}{thm:p03-w03-p03-032-spectrum}{Stability spectrum}
\begin{proof}
\emph{The bound for stable \(T\).}  If no formula has the order property,
then for each formula \(\varphi(x;y)\) with \(x\) a single variable,
\cref{thm:p03-w03-p03-032-formula} gives
\(\abs{S_\varphi(A)}\leq\abs A+\aleph_0\), and
\cref{thm:p03-w03-p03-032-product} gives
\(\abs{S_1(A)}\leq(\abs A+\aleph_0)^{\abs T}\).

(1)\(\Rightarrow\)(3).  Let \(\lambda^{\abs T}=\lambda\); then \(\lambda\) is
infinite, and for \(\abs A\leq\lambda\) the bound gives
\(\abs{S_1(A)}\leq\lambda^{\abs T}=\lambda\).

(3)\(\Rightarrow\)(2).  The cardinal \(\lambda=2^{\abs T}\) is infinite and
satisfies \(\lambda^{\abs T}=2^{\abs T\cdot\abs T}=\lambda\).

(2)\(\Rightarrow\)(1), together with the final claim.  Suppose
\(\varphi(x;y)\) has the order property, with \(x\) of length \(m\), and let
\(\lambda\) be any infinite cardinal.  By
\cref{thm:p03-w03-p03-032-formula} there is \(A\) with \(\abs A\leq\lambda\)
and \(\abs{S_\varphi(A)}>\lambda\).  Each \(r\in S_\varphi(A)\) is
consistent and extends to some \(r^*\in S_m(A)\); distinct \(r\) give
distinct \(r^*\), because \(r=r^*\mathbin{\upharpoonright}\varphi\).  So
\(\abs{S_m(A)}>\lambda\), and by \cref{thm:p03-w03-p03-032-one-to-many}
\(T\) is not \(\lambda\)-stable.  Thus an unstable \(T\) is
\(\lambda\)-stable for no infinite \(\lambda\), which is the contrapositive
of (2)\(\Rightarrow\)(1).
\end{proof}
\end{annexitem}

\begin{annexitem}{annex:p03-w03-p03-032-degree-basics}{thm:p03-w03-p03-032-degree-basics}{Basic properties of the degree}
\begin{proof}
Write \(\kappa_T=\abs T^+\) as in \cref{def:p03-w03-p03-032-degree}.

(1) By induction on \(\alpha\) we show: if
\(\mathfrak C\entails\forall x\,(\theta\to\theta')\) and
\(\operatorname{Deg}\theta\geq\alpha\), then
\(\operatorname{Deg}\theta'\geq\alpha\), for all formulas
\(\theta,\theta'\).  For \(\alpha=0\), a realization of \(\theta\) realizes
\(\theta'\).  For a limit \(\alpha\) the claim for all smaller ordinals gives
it.  For \(\alpha=\gamma+1\), take \(\psi\), \(k\) and \((b_l:l<\kappa_T)\)
witnessing \(\operatorname{Deg}\theta\geq\gamma+1\).  Each
\(\theta\wedge\psi(x;b_l)\) implies \(\theta'\wedge\psi(x;b_l)\), so by the
induction hypothesis \(\operatorname{Deg}(\theta'\wedge\psi(x;b_l))\geq\gamma\),
and the same \(\psi\), \(k\), \((b_l)\) witness
\(\operatorname{Deg}\theta'\geq\gamma+1\).  Taking \(\alpha\) arbitrary, the
set of ordinals \(\alpha\) with \(\operatorname{Deg}\theta\geq\alpha\) is
contained in the corresponding set for \(\theta'\); by (2) below these sets
determine the degrees, so \(\operatorname{Deg}\theta\leq\operatorname{Deg}\theta'\).

(2) By induction on \(\beta\) we show: \(\operatorname{Deg}\theta\geq\beta\)
implies \(\operatorname{Deg}\theta\geq\alpha\) for all \(\alpha\leq\beta\).
For \(\beta=0\) there is nothing to show.  For a limit \(\beta\) the
definition gives every \(\alpha<\beta\).  For \(\beta=\gamma+1\), a witness
\(\psi\), \(k\), \((b_l)\) gives
\(\operatorname{Deg}(\theta\wedge\psi(x;b_0))\geq\gamma\); the formula
\(\theta\wedge\psi(x;b_0)\) implies \(\theta\), so the induction on (1) at
the ordinal \(\gamma\), which used only the definition, gives
\(\operatorname{Deg}\theta\geq\gamma\), and the induction hypothesis for
\(\gamma\) gives every \(\alpha\leq\gamma\).

Consequently \(S_\theta\defeq\{\alpha:\operatorname{Deg}\theta\geq\alpha\}\)
is an initial segment of the ordinals.  If it is empty, \(\theta\) is not
realized and \(\operatorname{Deg}\theta=-1\).  If it is a proper nonempty
initial segment, it has a largest element: otherwise its supremum
\(\delta\) is a limit ordinal with \(\operatorname{Deg}\theta\geq\alpha\)
for all \(\alpha<\delta\), so \(\delta\in S_\theta\) by clause~(2) of the
definition, and \(\delta\) would be the largest element.  That largest
element is the unique \(\alpha\) with \(\operatorname{Deg}\theta\geq\alpha\)
and not \(\operatorname{Deg}\theta\geq\alpha+1\).  Otherwise \(S_\theta\)
contains all ordinals and \(\operatorname{Deg}\theta=\infty\).  In each case
\(\operatorname{Deg}\theta\) is well defined, and
\(\operatorname{Deg}\theta\leq\operatorname{Deg}\theta'\) holds exactly when
\(S_\theta\subseteq S_{\theta'}\), which completes (1).

(3) By induction on \(\alpha\) we show, for every formula \(\theta(x;y)\) of
\(\Lang\) and all tuples \(a\equiv_\varnothing a'\): if
\(\operatorname{Deg}\theta(x;a)\geq\alpha\) then
\(\operatorname{Deg}\theta(x;a')\geq\alpha\).  For \(\alpha=0\),
\(\exists x\,\theta(x;y)\) is a formula of \(\operatorname{tp}(a)\), hence of
\(\operatorname{tp}(a')\).  Limits are immediate.  For \(\alpha=\gamma+1\)
take witnesses \(\psi\), \(k\), \(e=(b_l:l<\kappa_T)\) for \(a\).  The tuple
\(e\) is small, so by \cref{thm:p03-w03-p03-031-saturation}\,(2) there is
\(e'=(b'_l:l<\kappa_T)\) with \(ae\equiv_\varnothing a'e'\).  For each \(w\)
with \(k\) elements the formula
\(\neg\exists x\bigwedge_{l\in w}\psi(x;y_l)\) holds of \((b_l)_{l\in w}\),
hence of \((b'_l)_{l\in w}\), so \((\psi(x;b'_l))\) is \(k\)-inconsistent.
For each \(l\), \(ab_l\equiv_\varnothing a'b'_l\), and the induction
hypothesis applied to the formula \(\theta(x;y)\wedge\psi(x;z)\) gives
\(\operatorname{Deg}(\theta(x;a')\wedge\psi(x;b'_l))\geq\gamma\).  So
\(\operatorname{Deg}\theta(x;a')\geq\gamma+1\).  By symmetry the two sets
\(S_{\theta(x;a)}\) and \(S_{\theta(x;a')}\) coincide.

(4) Each \(\theta_i\) implies \(\theta_1\vee\theta_2\), so (1) gives
\(\max(\operatorname{Deg}\theta_1,\operatorname{Deg}\theta_2)\leq
\operatorname{Deg}(\theta_1\vee\theta_2)\).  For the other inequality we show
by induction on \(\alpha\), for all pairs of formulas: if
\(\operatorname{Deg}(\theta_1\vee\theta_2)\geq\alpha\) then
\(\operatorname{Deg}\theta_i\geq\alpha\) for some \(i\in\{1,2\}\).  For
\(\alpha=0\), a realization of the disjunction realizes one disjunct.  For
\(\alpha=\gamma+1\), take witnesses \(\psi\), \(k\), \((b_l:l<\kappa_T)\).
Since \((\theta_1\vee\theta_2)\wedge\psi(x;b_l)\) is equivalent to
\((\theta_1\wedge\psi(x;b_l))\vee(\theta_2\wedge\psi(x;b_l))\), (1) and the
induction hypothesis give, for each \(l\), some \(i(l)\in\{1,2\}\) with
\(\operatorname{Deg}(\theta_{i(l)}\wedge\psi(x;b_l))\geq\gamma\).  As
\(\kappa_T\) is infinite, one value \(i\) is taken on a set
\(L\subseteq\kappa_T\) of cardinality \(\kappa_T\).  Re-indexing
\((b_l:l\in L)\) by \(\kappa_T\) gives a family that is still
\(k\)-inconsistent, every member of which has degree at least \(\gamma\)
together with \(\theta_i\); so \(\operatorname{Deg}\theta_i\geq\gamma+1\).
For a limit \(\alpha\), for each \(\beta<\alpha\) some \(i(\beta)\) has
\(\operatorname{Deg}\theta_{i(\beta)}\geq\beta\); one value \(i\) occurs for
cofinally many \(\beta<\alpha\), and by (2)
\(\operatorname{Deg}\theta_i\geq\beta\) for every \(\beta<\alpha\), hence
\(\operatorname{Deg}\theta_i\geq\alpha\).

(5) By (3), \(\operatorname{Deg}\theta(x;a)\) depends only on the formula
\(\theta(x;y)\) of \(\Lang\) and on the complete type \(\operatorname{tp}(a)\)
of the finite tuple \(a\) over \(\varnothing\).  There are at most
\(\abs T\) formulas and at most \(2^{\abs T}\) such types, so the set \(V\)
of ordinals that occur as degrees of formulas in \(x\) has cardinality at
most \(2^{\abs T}\).  Let \(\alpha_0\) be the least ordinal not in \(V\);
then \(\alpha_0<(2^{\abs T})^+\).  We show by induction on
\(\beta\geq\alpha_0\): for every formula \(\theta\) in \(x\),
\(\operatorname{Deg}\theta\geq\alpha_0\) implies
\(\operatorname{Deg}\theta\geq\beta\).  The case \(\beta=\alpha_0\) is the
hypothesis, and limits follow from the definition.  Let
\(\operatorname{Deg}\theta\geq\alpha_0\) and let the claim hold for
\(\beta\).  The degree of \(\theta\) is \(\infty\) or an ordinal in \(V\) that
is at least \(\alpha_0\), hence different from \(\alpha_0\); in both cases
\(\operatorname{Deg}\theta\geq\alpha_0+1\).  Take witnesses \(\psi\), \(k\),
\((b_l)\): each \(\theta\wedge\psi(x;b_l)\) has degree at least
\(\alpha_0\), hence at least \(\beta\) by the induction hypothesis, and so
\(\operatorname{Deg}\theta\geq\beta+1\).  Thus
\(\operatorname{Deg}\theta\geq\beta\) for all \(\beta\), that is,
\(\operatorname{Deg}\theta=\infty\).
\end{proof}
\end{annexitem}

\begin{annexitem}{annex:p03-w03-p03-032-dividing-degree}{thm:p03-w03-p03-032-dividing-degree}{Dividing lowers the degree}
\begin{proof}
Let \(f\) be the tuple of parameters of \(\theta\), so \(\theta=\theta(x;f)\)
with \(f\) from \(F\), and let \(\alpha=\operatorname{Deg}\theta\).  By
\cref{def:p03-dividing-forking} and
\cref{thm:p03-w03-p03-031-saturation}\,(3) there is an \(F\)-indiscernible
sequence \((e_i:i<\omega)\) in \(\mathfrak C\) with \(e_0=e\) such that
\(\{\varphi(x;e_i):i<\omega\}\) is not realized.  By realization by
saturation, some finite subset is not realized; list its indices as
\(i_0<\cdots<i_{k-1}\).  The formula
\(\neg\exists x\bigwedge_{r<k}\varphi(x;y_r)\) holds of
\((e_{i_0},\ldots,e_{i_{k-1}})\), so by indiscernibility it holds of every
increasing \(k\)-tuple from the sequence; as every \(k\)-element set of
indices can be listed increasingly, \((\varphi(x;e_i):i<\omega)\) is
\(k\)-inconsistent.

Let \(\Gamma\) be the set of formulas, in variables \(z_l\) for
\(l<\kappa_T\) of the length of \(e\), consisting of
\(\rho(z_{l_0},\ldots,z_{l_{n-1}})\) for every formula \(\rho\) with
parameters from \(F\), every \(n\) and every \(l_0<\cdots<l_{n-1}<\kappa_T\)
such that \(\mathfrak C\entails\rho(e_0,\ldots,e_{n-1})\).  A finite subset of
\(\Gamma\) mentions finitely many indices; sending them in increasing order
to \(0,1,2,\ldots\) and the variables to the corresponding \(e_i\) satisfies
it, by indiscernibility.  The number of variables is \(\kappa_T\), which is
small, so \(\Gamma\) is realized by some \((e'_l:l<\kappa_T)\).  Each
\(e'_l\) satisfies every formula over \(F\) true of \(e_0\), so
\(e'_l\equiv_F e\), and every increasing \(k\)-tuple of the \(e'_l\)
satisfies \(\neg\exists x\bigwedge_{r<k}\varphi(x;y_r)\); hence
\((\varphi(x;e'_l):l<\kappa_T)\) is \(k\)-inconsistent.

Since \(f\) is from \(F\), \(fe'_l\equiv_\varnothing fe\) for every \(l\), and
\cref{thm:p03-w03-p03-032-degree-basics}\,(3), applied to the formula
\(\theta(x;w)\wedge\varphi(x;z)\), gives
\(\operatorname{Deg}(\theta\wedge\varphi(x;e'_l))=\beta\) for all \(l\),
where \(\beta\defeq\operatorname{Deg}(\theta\wedge\varphi(x;e))\).  If
\(\beta\geq\alpha\), then \(\varphi\), \(k\) and \((e'_l)\) witness
\(\operatorname{Deg}\theta\geq\alpha+1\), which is false because \(\alpha\)
is an ordinal equal to \(\operatorname{Deg}\theta\).  Hence \(\beta<\alpha\).
\end{proof}
\end{annexitem}

\begin{annexitem}{annex:p03-w03-p03-032-degree-tree}{thm:p03-w03-p03-032-degree-tree}{Infinite degree gives inconsistency trees}
\begin{proof}
Write \(\kappa=\kappa_T\) and, for \(\nu\in\kappa^{\leq n}\) and
\(j\leq\abs\nu\), \(\nu\mathbin{\upharpoonright}j\) for the initial segment of
length \(j\).  Let \(\alpha_0\) be as in
\cref{thm:p03-w03-p03-032-degree-basics}\,(5).

\emph{Step 1: finite trees.}  We choose, by recursion on \(n\), formulas
\(\psi_n(x;y_n)\) and integers \(k_n\) and, for each \(n<\omega\), tuples
\(a^n_\nu\) for \(\nu\in\kappa^{\leq n}\setminus\{\varnothing\}\), such that
\(\psi_1,k_1,\ldots,\psi_n,k_n\) are fixed at stage \(n\) and:
\begin{enumerate}
  \item[(a)] for every \(\eta\in\kappa^n\) the formula
        \(\theta^n_\eta\defeq\bigwedge_{1\leq j\leq n}
        \psi_j(x;a^n_{\eta\mathbin{\upharpoonright}j})\) has degree
        \(\infty\) (for \(n=0\) it is \(x=x\));
  \item[(b)] for every \(j<n\) and \(\nu\in\kappa^j\) the family
        \((\psi_{j+1}(x;a^n_{\nu\frown\langle i\rangle}):i<\kappa)\) is
        \(k_{j+1}\)-inconsistent.
\end{enumerate}
At \(n=0\) the family is empty and (a) is the hypothesis.  Assume stage
\(n\) is done.  For \(\eta\in\kappa^n\),
\(\operatorname{Deg}\theta^n_\eta=\infty\geq\alpha_0+1\), so there are
\(\psi^\eta\), \(k^\eta\) and \((b^\eta_l:l<\kappa)\), \(k^\eta\)-inconsistent,
with \(\operatorname{Deg}(\theta^n_\eta\wedge\psi^\eta(x;b^\eta_l))\geq
\alpha_0\), hence \(=\infty\) by
\cref{thm:p03-w03-p03-032-degree-basics}\,(5).  Put
\(\ell(\eta)=(\psi^\eta,k^\eta)\); these labels range over a set \(\Lambda\) of
cardinality at most \(\abs T<\kappa\).

Define \(\ell^*(\nu)\in\Lambda\) for \(\nu\in\kappa^{\leq n}\) by downward
recursion on \(\abs\nu\): \(\ell^*(\eta)=\ell(\eta)\) when \(\abs\eta=n\),
and for \(\abs\nu<n\) let \(\ell^*(\nu)\) be a label \(\ell\) such that
\(\{i<\kappa:\ell^*(\nu\frown\langle i\rangle)=\ell\}\) has cardinality
\(\kappa\).  Such a label exists because \(\kappa\) is regular and greater
than \(\abs\Lambda\), so \(\kappa\) is not a union of \(\abs\Lambda\) sets
each of cardinality below \(\kappa\).  Let \(K\) be the set of
\(\nu\in\kappa^{\leq n}\) with \(\ell^*(\nu\mathbin{\upharpoonright}j)=
\ell^*(\varnothing)\) for all \(j\leq\abs\nu\).  Then \(\varnothing\in K\),
\(K\) is closed under initial segments, and each \(\nu\in K\) with
\(\abs\nu<n\) has \(\kappa\) children
\(\nu\frown\langle i\rangle\in K\), namely those with
\(\ell^*(\nu\frown\langle i\rangle)=\ell^*(\nu)\).  Choose
\(g:\kappa^{\leq n}\to K\) by recursion: \(g(\varnothing)=\varnothing\) and
\(g(\nu\frown\langle i\rangle)=g(\nu)\frown\langle h_\nu(i)\rangle\) for an
injection \(h_\nu\) of \(\kappa\) into
\(\{i':g(\nu)\frown\langle i'\rangle\in K\}\).  Then \(g\) preserves length
and initial segments and is injective on the children of each node.

Let \((\psi_{n+1},k_{n+1})\defeq\ell^*(\varnothing)\), and put
\(a^{n+1}_\nu\defeq a^n_{g(\nu)}\) for \(1\leq\abs\nu\leq n\) and
\(a^{n+1}_{\eta\frown\langle l\rangle}\defeq b^{g(\eta)}_l\) for
\(\eta\in\kappa^n\), \(l<\kappa\).  For \(\eta\in\kappa^n\), \(g(\eta)\in K\)
has length \(n\), so \(\ell(g(\eta))=\ell^*(\varnothing)\), that is,
\(\psi^{g(\eta)}=\psi_{n+1}\) and \(k^{g(\eta)}=k_{n+1}\).  Hence
\(\theta^{n+1}_{\eta\frown\langle l\rangle}=\theta^n_{g(\eta)}\wedge
\psi_{n+1}(x;b^{g(\eta)}_l)\), which has degree \(\infty\); this is (a).  For
(b) at a level \(j<n\), the family at \(\nu\) is the subfamily of the
\(k_{j+1}\)-inconsistent family at \(g(\nu)\) indexed by the injective
\(h_\nu\), so it is \(k_{j+1}\)-inconsistent; at level \(n\) it is the family
\((\psi_{n+1}(x;b^{g(\eta)}_l):l<\kappa)\), which is
\(k_{n+1}\)-inconsistent.  Finally \(k_{n+1}\geq2\): each
\(\psi_{n+1}(x;b^{g(\eta)}_l)\) is realized, since
\(\theta^{n+1}_{\eta\frown\langle l\rangle}\) has degree at least \(0\), so a
\(1\)-inconsistent family is impossible.

\emph{Step 2: compactness.}  Let \(\lambda\) be infinite and small, and take
variables \(z_\eta\) for \(\eta\in\lambda^{<\omega}\setminus\{\varnothing\}\),
of the length of \(y_{\abs\eta}\).  Let \(\Sigma\) consist of the formulas
\[
  \exists x\bigwedge_{1\leq j\leq n}\psi_j(x;z_{\nu\mathbin{\upharpoonright}j})
  \qquad(n\geq1,\ \nu\in\lambda^n)
\]
and
\[
  \neg\exists x\bigwedge_{r<k_{n+1}}
  \psi_{n+1}(x;z_{\eta\frown\langle i_r\rangle})
  \qquad(n<\omega,\ \eta\in\lambda^n,\
  i_0<\cdots<i_{k_{n+1}-1}<\lambda).
\]
A finite \(\Sigma_0\subseteq\Sigma\) mentions finitely many nodes; let \(N\)
be the finite set of these nodes and their initial segments and \(n\) the
greatest length in \(N\).  Define \(f:N\to\kappa^{\leq n}\) by
\(f(\varnothing)=\varnothing\) and by sending the finitely many children of
\(\nu\) in \(N\) injectively to children of \(f(\nu)\); this preserves
length and initial segments.  Interpret \(z_\eta\) as \(a^n_{f(\eta)}\).  A
formula of the first kind, for \(\nu\in N\) of length \(n'\leq n\), holds
because \(\theta^n_{\eta}\) is realized for any \(\eta\in\kappa^n\)
extending \(f(\nu)\) and implies
\(\bigwedge_{j\leq n'}\psi_j(x;a^n_{f(\nu\mathbin{\upharpoonright}j)})\).  A
formula of the second kind holds because distinct children are sent to
distinct children and the family at \(f(\eta)\) is
\(k_{n'+1}\)-inconsistent by (b), where \(n'+1\leq n\).  So \(\Sigma_0\) is
satisfied.  \(\Sigma\) has no parameters and \(\lambda\) times finitely many
variables, so it is realized by saturation; let \((a_\eta)\) realize it.
Then (2) holds by the second kind of formula.  For (1), every finite subset
of \(\{\psi_n(x;a_{\nu\mathbin{\upharpoonright}n}):n\geq1\}\) is contained in
\(\{\psi_j(x;a_{\nu\mathbin{\upharpoonright}j}):1\leq j\leq n\}\) for some
\(n\), which is realized by the first kind of formula; by realization by
saturation the whole set is realized.
\end{proof}
\end{annexitem}

\begin{annexitem}{annex:p03-w03-p03-032-dividing-tree-count}{thm:p03-w03-p03-032-dividing-tree-count}{Inconsistency trees give many types}
\begin{proof}
There are \(\abs{\lambda^{<\omega}}=\lambda\) nodes and each \(a_\eta\) is
finite, so \(\abs A\leq\lambda\).  For \(\nu\in\lambda^\omega\) choose
\(c_\nu\) realizing \(\{\psi_n(x;a_{\nu\mathbin{\upharpoonright}n}):n\geq1\}\)
and let \(q_\nu=\operatorname{tp}(c_\nu/A)\in S_m(A)\).

Fix \(q\in S_m(A)\) and let \(N_q=\{\nu:q_\nu=q\}\).  For a node \(\eta\),
let \(C_q(\eta)\) be the set of \(i<\lambda\) with
\(\psi_{\abs\eta+1}(x;a_{\eta\frown\langle i\rangle})\in q\).  It has fewer
than \(k_{\abs\eta+1}\) elements: otherwise \(q\), which is realized,
contains a \(k_{\abs\eta+1}\)-element subfamily of a
\(k_{\abs\eta+1}\)-inconsistent family.  If \(\nu\in N_q\), then
\(\nu(n)\in C_q(\nu\mathbin{\upharpoonright}n)\) for every \(n\).  Send
\(\nu\in N_q\) to the sequence \((j_n)_{n<\omega}\) where \(\nu(n)\) is the
\(j_n\)-th element of \(C_q(\nu\mathbin{\upharpoonright}n)\) in increasing
order.  The sequence determines \(\nu\) by recursion on \(n\), and
\(j_n<k_{n+1}\), so \(\abs{N_q}\leq\abs{\prod_{n<\omega}k_{n+1}}
\leq\aleph_0^{\aleph_0}=2^{\aleph_0}\).

The sets \(N_q\) partition \(\lambda^\omega\), which has cardinality
\(\lambda^{\aleph_0}\), so
\(\lambda^{\aleph_0}\leq\abs{S_m(A)}\cdot2^{\aleph_0}\).  If
\(\abs{S_m(A)}\leq\lambda\), then, as \(\lambda\geq2^{\aleph_0}\),
\(\lambda^{\aleph_0}\leq\lambda\cdot2^{\aleph_0}=\lambda\), contrary to the
hypothesis.  So \(\abs{S_m(A)}>\lambda\) with \(\abs A\leq\lambda\).  If
\(T\) were \(\lambda\)-stable, \cref{thm:p03-w03-p03-032-one-to-many} would
give \(\abs{S_m(A)}\leq\lambda\); hence \(T\) is not \(\lambda\)-stable.
\end{proof}
\end{annexitem}

\begin{annexitem}{annex:p03-w03-p03-032-pressing-down}{thm:p03-w03-p03-032-pressing-down}{A bounded regressive function on many limit ordinals}
\begin{proof}
For \(j<\lambda\) put
\(S_j=\{k<\lambda:k\text{ is a nonzero limit ordinal},\ f(k)\leq j\}\).
Suppose every \(S_j\) has cardinality below \(\lambda\).  Regularity makes
each \(S_j\) bounded in \(\lambda\).  Choose a strictly increasing sequence
\((a_n:n<\omega)\) below \(\lambda\) with
\(a_{n+1}>\sup S_{a_n}\), and put \(\delta=\sup_n a_n\).
Uncountable regularity gives \(\delta<\lambda\), and \(\delta\) is a
nonzero limit ordinal.  Since \(f(\delta)<\delta\), some \(n\) has
\(f(\delta)\leq a_n\).  Then \(\delta\in S_{a_n}\), contradicting
\(\delta\geq a_{n+1}>\sup S_{a_n}\).  Thus some \(S_j\) has cardinality
\(\lambda\), as asserted.
\end{proof}
\end{annexitem}

\begin{annexitem}{annex:p03-w03-p03-032-uniform}{thm:p03-w03-p03-032-uniform}{Finite uniformity on a sequence of regular uncountable length}
\begin{proof}
We induct on \(N\).  For \(N=1\), finitely many truth values partition
\(\lambda\) into finitely many classes, one of size \(\lambda\).  Suppose
\(N>1\).  Put \(A_i=D\cup\bigcup_{j<i}\operatorname{rng}(b_j)\).
For \(\delta\in\Delta\), regard its last tuple as object variables:
\(\varphi_\delta(z;y_0,\ldots,y_{N-2})=
\delta(y_0,\ldots,y_{N-2},z)\).  The fixed parameters from \(D\) may be
named here.  Each of these formulas is stable: an order pattern for it
would be an order pattern for the corresponding formula of \(\Lang\)
with those parameters included in its parameter tuple.  Hence its local
tree rank is finite by \cref{thm:p03-w03-p03-032-spectrum}.

For every nonzero limit \(i<\lambda\), choose a finite part
\(q_{\delta,i}\) of \(\operatorname{tp}_{\varphi_\delta}(b_i/A_i)\)
of the same rank, using \cref{thm:p03-w03-p03-032-rank-basics}.  Finitely
many such finite parts use, outside \(D\), entries of only finitely many
\(b_j\) with \(j<i\).  Their indices are bounded by some \(f(i)<i\).
By \cref{thm:p03-w03-p03-032-pressing-down}, on a set of \(\lambda\)
indices these supports lie in one fixed \(A_{j+1}\).  Discard the bounded
initial segment at or below \(j+1\).  There are fewer than \(\lambda\)
possible tuples of finite sets \((q_{\delta,i}:\delta\in\Delta)\) over
\(A_{j+1}\): each uses just a fixed formula and finitely many parameters,
and \(\max(\abs{A_{j+1}},\aleph_0)<\lambda\).  Regularity therefore gives
a set \(I\) of size \(\lambda\) on which these finite sets are fixed;
write them as \(q_\delta\), with ranks \(r_\delta\).

For a parameter tuple \(a\) from \(A_i\), and any \(t\in I\) with
\(t\geq i\), \cref{thm:p03-w03-p03-032-local-character} says
\[
 \varphi_\delta(b_t;a)\quad\Longleftrightarrow\quad
 \mathrm R_{\varphi_\delta}
       (q_\delta\cup\{\varphi_\delta(z;a)\})\geq r_\delta.
\]
Indeed \(q_\delta\) is a finite part of exactly the rank of the complete
local type of \(b_t\) over \(A_t\).  Thus the truth value on the left is
independent of every sufficiently late \(t\in I\).  For all \(\delta\)
and all parameter tuples from \(D\cup\bigcup_{i<\lambda}
\operatorname{rng}(b_i)\), collect these eventual signed instances in the
variable \(z\).  Every finite subset is realized by a sufficiently late
\(b_t\) with \(t\in I\).  Saturation supplies one tuple \(b_*\)
realizing them all.

It follows that for \(k_0<\cdots<k_{N-1}\) in \(I\),
\[
 \delta(b_{k_0},\ldots,b_{k_{N-1}})
 \quad\Longleftrightarrow\quad
 \delta(b_{k_0},\ldots,b_{k_{N-2}},b_*).
\]
The set \(I\), in its inherited order, has order type \(\lambda\).
Apply the induction hypothesis to this sequence, the finite family of
\((N-1)\)-tuple formulas on the right, and the parameter set
\(D\cup\operatorname{rng}(b_*)\), still of size below \(\lambda\).
The resulting set \(K\subseteq I\) has size \(\lambda\) and has the
required uniformity for every \(\delta\in\Delta\).
\end{proof}
\end{annexitem}

\begin{annexitem}{annex:p03-w03-p03-032-same-degree}{thm:p03-w03-p03-032-same-degree}{Counting complete extensions that keep an ordinal degree}
\begin{proof}
If \(\theta\) is inconsistent, \(P_\theta\) is empty.  Otherwise put
\(\kappa=\abs T^+\), and suppose (1) fails for \(\varphi\).  Choose
\((p_i:i<\kappa)\) in \(P_\theta\) with pairwise distinct complete local
types \(s_i=p_i\mathbin{\upharpoonright}\varphi\).  Stability and
\cref{thm:p03-w03-p03-032-rank-basics} give finite
\(q_i\subseteq s_i\) with \(\mathrm R_\varphi(q_i)=
\mathrm R_\varphi(s_i)<\omega\).  There are only countably many choices
of the finite number of instances, their signs, and this rank.  Thin to
\(\kappa\) indices with all three fixed, and write
\(\bigwedge q_i=\psi(x;c_i)\), using one formula \(\psi\); call the
common local rank \(r\).

For \(i\neq j\), \(p_i\) cannot contain \(q_j\).  Otherwise both
\(s_i\) and \(s_j\) would be complete local extensions of \(q_j\) of
rank \(r\), contrary to the uniqueness in
\cref{thm:p03-w03-p03-032-local-character}.  Since \(p_i\) is complete,
it therefore contains \(\psi(x;c_i)\) and every
\(\neg\psi(x;c_j)\) with \(j\neq i\).  Consequently the conjunction of
\(\theta\), this positive instance, and any finitely many of those
negative instances has degree exactly \(\alpha\): it belongs to the
type \(p_i\) of degree \(\alpha\), and it implies \(\theta\).

There is an integer \(h\geq1\) such that no \(h\) parameter tuples for
\(\psi\) admit every signed pattern.  If there were no such integer,
the patterns coding initial segments would give every finite order pattern
for \(\psi\), and compactness would contradict stability.  Let \(\Delta\)
contain the finitely many formulas
\[
 \exists x\,\bigwedge_{j<2h}\psi(x;y_j)^{\varepsilon(j)}
 \qquad(\varepsilon\in\{0,1\}^{2h}).
\]
Apply \cref{thm:p03-w03-p03-032-uniform} to \((c_i:i<\kappa)\) and
\(\Delta\), and reindex the resulting sequence by \(\kappa\).  Put
\(\rho(x;u,v)=\neg\psi(x;u)\wedge\psi(x;v)\).  We claim that
\[
       \{\rho(x;c_{2i},c_{2i+1}):i<\kappa\}
\]
is \(h\)-inconsistent.  A realization of any \(h\) of its members would
give the alternating negative--positive pattern on their \(2h\) ordered
parameters.  Uniformity would give that pattern on every ordered
\(2h\)-tuple from the sequence.  This would shatter the fixed parameters
\(c_{3j+1}\), \(j<h\): for any desired signs choose, in the \(j\)-th
block, the ordered pair \((c_{3j},c_{3j+1})\) for a positive sign and
\((c_{3j+1},c_{3j+2})\) for a negative sign.  The alternating pattern on
these \(h\) pairs realizes exactly the desired signs on the fixed
parameters.  This contradicts the choice of \(h\).

On the other hand, each
\(\theta\wedge\rho(x;c_{2i},c_{2i+1})\) belongs to \(p_{2i+1}\), so
has degree \(\alpha\).  The displayed \(h\)-inconsistent family has
\(\kappa_T\) members and therefore witnesses
\(\operatorname{Deg}\theta\geq\alpha+1\), contradicting its degree
\(\alpha\).  This proves (1).  For (2), every complete type is determined
by its restrictions to the at most \(\abs T\) partitioned formulas.
There are at most \(\abs T\) choices for each restriction by (1), whence
\(\abs{P_\theta}\leq\abs T^{\abs T}=2^{\abs T}\).
\end{proof}
\end{annexitem}

\begin{annexitem}{annex:p03-w03-p03-032-forking-chain}{thm:p03-w03-p03-032-forking-chain}{From an inconsistency tree to failure of finite local character}
\begin{proof}
Take the formulas \(\psi_n\), integers \(k_n\) and an
\(\omega\)-branching tree from \cref{thm:p03-w03-p03-032-degree-tree}.
We construct rows \((d_{n,i}:i<\omega)\), \(n\geq1\), such that:
\begin{enumerate}
  \item the row numbered \(n\) is indiscernible over all entries of
        \(d_{1,0},\ldots,d_{n-1,0}\);
  \item \(\{\psi_n(x;d_{n,i}):i<\omega\}\) is \(k_n\)-inconsistent;
  \item \(\{\psi_n(x;d_{n,0}):n\geq1\}\) is consistent.
\end{enumerate}
Each requirement is expressed by a set of first-order formulas in the row
variables: equality of truth values for increasing tuples in (1), the
negated existential conjunction for each \(k_n\)-tuple in (2), and the
existential assertion for each finite initial conjunction in (3).

To verify finite satisfiability, fix a finite set of these requirements
and a height \(N\) containing its row indices.  Starting at the root of
the original tree, proceed down one branch.  At stage \(n\), the earlier
chosen branch parameters are fixed.  Only finitely many formulas over
those parameters occur in the required indiscernibility conditions for
row \(n\).  Repeated applications of the infinite Ramsey theorem
(\cref{thm:p03-w03-p03-028-ramsey-a}) give an infinite subsequence of the
current node's children homogeneous for all those formulas.  Assign the
finitely many mentioned row positions, including position zero, to
increasing members of this subsequence; choose the member at position
zero as the next node on the branch.  Continue inside that node's full
subtree.  Sibling inconsistency survives thinning, and all chosen
zero-position parameters lie on a branch of the original tree, so their
finite conjunction is realized.  This satisfies the fixed finite set of
requirements.  Compactness and saturation give all the rows at once.
This is the same finite-Ramsey argument as
\cref{thm:p03-w03-p03-035-em-extraction}, with the earlier branch
parameters named separately at each row.

Let \(B\) be the union of the entries of the \(d_{n,0}\), choose \(a\)
realizing (3), and set \(p=\operatorname{tp}(a/B)\).  The set \(B\) is
countable.  Conditions (1) and (2), with the first element \(d_{n,0}\),
witness that \(\psi_n(x;d_{n,0})\) divides over
\(B_{<n}=\bigcup_{j<n}\operatorname{rng}(d_{j,0})\).  Every finite
\(F\subseteq B\) is contained in some \(B_{<n}\).  The same row is
\(F\)-indiscernible, so its first formula divides over \(F\).  That
formula belongs to \(p\), which therefore forks over \(F\), as required.
\end{proof}
\end{annexitem}


\begin{annexitem}{annex:p03-w03-p03-032-superstable}{thm:p03-w03-p03-032-superstable}{Superstability above \(2^{\abs T}\)}
\begin{proof}
\emph{(1) implies (2).}  Put \(\mu_0=2^{\abs T}\), \(\mu_{n+1}=2^{\mu_n}\)
and \(\lambda=\sup_{n<\omega}\mu_n\).  By Cantor's theorem
\(\mu_n<\mu_{n+1}\leq\lambda\), so \(\lambda\) is the supremum of a
countable sequence of smaller cardinals, and
\(\lambda^{\aleph_0}>\lambda\) by \cref{thm:p03-w03-p03-032-koenig}.
Moreover \(\lambda\geq2^{\abs T}\geq2^{\aleph_0}\), and \(T\) is
\(\lambda\)-stable by (1).

\emph{(2) implies (3).}  Let \(T\) be \(\lambda\)-stable with
\(\lambda^{\aleph_0}>\lambda\geq2^{\aleph_0}\).  Then \(T\) is stable by
\cref{thm:p03-w03-p03-032-spectrum}, since it is \(\lambda\)-stable for an
infinite \(\lambda\).  If \(\operatorname{Deg}_m(x=x)=\infty\) for some
\(m\), \cref{thm:p03-w03-p03-032-degree-tree} gives a tree of width
\(\lambda\) and \cref{thm:p03-w03-p03-032-dividing-tree-count} shows that
\(T\) is not \(\lambda\)-stable, a contradiction.  So
\(\operatorname{Deg}_m(x=x)<\infty\) for every \(m\).

\emph{(3) implies (4).}  Let \(T\) satisfy (3), and let \(p\) be a complete
type in \(m\) variables over a set \(B\).  For every finite
\(\pi_0\subseteq p\), \(\bigwedge\pi_0\) implies \(x=x\), so
\(\operatorname{Deg}\bigwedge\pi_0\leq\operatorname{Deg}_m(x=x)<\infty\) by
\cref{thm:p03-w03-p03-032-degree-basics}\,(1).  The degrees of the finite
subsets of \(p\) are therefore ordinals, and the least of them is attained
by some finite \(q\subseteq p\): \(\operatorname{Deg}p=
\operatorname{Deg}\theta\), where \(\theta=\bigwedge q\); call this ordinal
\(\alpha\).  It
suffices to prove the final assertion for any such \(q\).  Let \(F\subseteq
B\) be the finite set of parameters of \(\theta\).

Suppose that \(p\) forks over \(F\).  By \cref{def:p03-dividing-forking} and
\cref{thm:p03-w03-p03-031-saturation}\,(3), \(p\) contains a formula
\(\chi\) and there are formulas \(\varphi_1(x;e_1),\ldots,\varphi_r(x;e_r)\)
with parameters in \(\mathfrak C\), each dividing over \(F\), such that
\(\mathfrak C\entails\forall x\,(\chi\to\bigvee_{j=1}^r\varphi_j(x;e_j))\).
Then \(\theta\wedge\chi\) implies
\(\bigvee_{j=1}^r(\theta\wedge\varphi_j(x;e_j))\), so by
\cref{thm:p03-w03-p03-032-degree-basics}\,(1) and~(4), the latter applied
\(r-1\) times,
\[
  \operatorname{Deg}(\theta\wedge\chi)\leq
  \max_{1\leq j\leq r}\operatorname{Deg}(\theta\wedge\varphi_j(x;e_j)).
\]
Each \(\varphi_j(x;e_j)\) divides over \(F\), \(\theta\) has parameters from
\(F\) and \(\operatorname{Deg}\theta=\alpha<\infty\), so
\cref{thm:p03-w03-p03-032-dividing-degree} gives
\(\operatorname{Deg}(\theta\wedge\varphi_j(x;e_j))<\alpha\) for every
\(j\), and hence \(\operatorname{Deg}(\theta\wedge\chi)<\alpha\).  But
\(\theta\wedge\chi=\bigwedge(q\cup\{\chi\})\) with \(q\cup\{\chi\}\) a finite
subset of \(p\), so \(\operatorname{Deg}(\theta\wedge\chi)\geq
\operatorname{Deg}p=\alpha\).  This contradiction shows that \(p\) does not
fork over the finite set \(F\subseteq B\).  As \(T\) is stable, it is
superstable in the sense of \cref{def:p03-stability-spectrum}.

\emph{(4) implies (3).}  Stability is part of the definition in (4).
If \(\operatorname{Deg}_m(x=x)=\infty\) for some \(m\),
\cref{thm:p03-w03-p03-032-forking-chain} supplies a complete type over a
countable set that forks over every finite subset, contrary to the finite
local character required in that definition.  Hence every such degree is
an ordinal.

\emph{(3) implies (5).}  Take \(m=1\).

\emph{(5) implies (1), and the type bound.}  Under (5), every formula in
one object variable has ordinal degree, by monotonicity.  For each
\(p\in S_1(A)\), the set of degrees of its finite conjunctions therefore
has a least member; choose a finite \(q_p\subseteq p\) attaining it and
put \(\theta_p=\bigwedge q_p\).  There are at most
\(\abs A+\abs T\) possible finite conjunctions over \(A\).  For a fixed
one, \cref{thm:p03-w03-p03-032-same-degree}\,(2) bounds the number of
complete extensions of its degree by \(2^{\abs T}\).  Thus
\[
 \abs{S_1(A)}\leq(\abs A+\abs T)\,2^{\abs T}
                   =\abs A+2^{\abs T}.
\]
If \(\abs A\leq\lambda\) and \(\lambda\geq2^{\abs T}\), this is at
most \(\lambda\), proving (1).  The implication (3) to (4) already proved
the additional nonforking assertion for every finite part of the same
degree.  All five conditions and both final assertions now follow.

\end{proof}
\end{annexitem}

\begin{annexitem}{annex:p03-w03-p03-032-qf-test}{thm:p03-w03-p03-032-qf-test}{Quantifier-free types decide types}
\begin{proof}
(1) We show, by induction on the formula \(\theta(\bar x)\), that for all
\(\omega\)-saturated \(\mathcal M,\mathcal N\entails T_0\), every finite
partial isomorphism \(f\) and every tuple \(\bar a\) from its domain,
\(\mathcal M\entails\theta(\bar a)\) iff \(\mathcal N\entails\theta(f\bar a)\).
For atomic \(\theta\) this is the definition, and the steps for \(\neg\) and
\(\wedge\) are immediate.  Let \(\theta=\exists y\,\chi(\bar x,y)\).  If
\(\mathcal M\entails\chi(\bar a,c)\), extend \(f\) to a finite partial
isomorphism \(g\) with \(c\) in its domain; by the induction hypothesis for
\(g\), \(\mathcal N\entails\chi(f\bar a,g(c))\), so
\(\mathcal N\entails\theta(f\bar a)\).  The converse uses an extension with a
given \(d\in N\) in its range and the induction hypothesis for its inverse,
which is a finite partial isomorphism from \(\mathcal N\) to \(\mathcal M\).

\emph{\(\omega\)-saturated elementary extensions.}  We show that every
structure \(\mathcal M\) has an \(\omega\)-saturated elementary extension.
First, \(\mathcal M\) has an elementary extension \(\mathcal M'\) realizing
every set \(p(x)\) of formulas in one variable with parameters from a finite
subset of \(M\) each of whose finite subsets is satisfied in \(\mathcal M\).
Indeed, let \(\Sigma\) consist of the elementary diagram of \(\mathcal M\)
(with a constant \(c_e\) for each \(e\in M\)) and, for each such \(p\), the
formulas \(p(d_p)\) for a new constant \(d_p\).  A finite subset of
\(\Sigma\) involves finitely many of the sets \(p\), and for each only a
finite part, satisfied in \(\mathcal M\) by some element; interpreting
\(c_e\) as \(e\) and \(d_p\) as that element satisfies it in \(\mathcal M\).
By first-order compactness (\cref{thm:fol-compactness}) \(\Sigma\) has a
model; since it satisfies the elementary diagram, \(e\mapsto c_e\) is an
elementary embedding of \(\mathcal M\), and identifying \(\mathcal M\) with
its image gives \(\mathcal M'\).  Now put \(\mathcal M_0=\mathcal M\),
\(\mathcal M_{n+1}=\mathcal M_n'\) and let \(\mathcal M^*\) be the union of
the chain.  By induction on formulas, \(\mathcal M_n\preccurlyeq\mathcal
M^*\) for every \(n\): the only nontrivial step is that
\(\mathcal M^*\entails\exists y\,\chi(\bar a,y)\) with \(\bar a\) from
\(M_n\) and a witness \(e\in M_k\), \(k\geq n\), gives
\(\mathcal M_k\entails\chi(\bar a,e)\) by the induction hypothesis for
\(\mathcal M_k\), hence \(\mathcal M_k\entails\exists y\,\chi(\bar a,y)\), and
then \(\mathcal M_n\entails\exists y\,\chi(\bar a,y)\) because
\(\mathcal M_n\preccurlyeq\mathcal M_k\).  If \(p\) has parameters from a
finite \(F\subseteq M^*\) and each finite part of \(p\) is satisfied in
\(\mathcal M^*\), then \(F\subseteq M_n\) for some \(n\), each finite part is
satisfied in \(\mathcal M_n\) by elementarity, \(p\) is realized in
\(\mathcal M_{n+1}\), and by elementarity in \(\mathcal M^*\).  So
\(\mathcal M^*\) is \(\omega\)-saturated.

(2) Let \(\mathcal M,\mathcal N\entails T_0\), and let
\(\mathcal M^*,\mathcal N^*\) be \(\omega\)-saturated elementary extensions.  The language has no constant
symbols, so the empty map is a finite partial isomorphism, and by (1) it
preserves every sentence: \(\mathcal M^*\equiv\mathcal N^*\), hence
\(\mathcal M\equiv\mathcal N\).  As \(T_0\) is consistent, it is complete.

(3) Let \(\bar a,\bar b\) satisfy the same atomic formulas in
\(\mathcal M\entails T_0\), and let \(\mathcal M^*\) be an
\(\omega\)-saturated elementary extension.  The map \(\bar a\mapsto\bar b\)
is well defined and injective, since equality is atomic, and it is a finite
partial isomorphism from \(\mathcal M^*\) to itself.  By (1) it preserves
every formula in \(\mathcal M^*\), hence in \(\mathcal M\).  For the second
sentence, a formula \(\theta(x;\bar a)\) over \(A\) holds of \(c\) iff
\(\theta(x;\bar y)\) holds of \((c,\bar a)\), and by the first sentence this
holds iff it holds of \((c',\bar a)\).
\end{proof}
\end{annexitem}

\begin{annexitem}{annex:p03-w03-p03-032-tp}{thm:p03-w03-p03-032-tp}{\(T_P\) is superstable and not \(\omega\)-stable}
\begin{proof}
(1) \emph{A model.}  On \(2^\omega\times\omega\) let \(P_n\) hold of
\((\tau,j)\) iff \(\tau(n)=1\).  For finite \(\sigma:I\to\{0,1\}\), every
\((\tau,j)\) with \(\tau\) extending \(\sigma\) satisfies \(P^\sigma\), and
there are infinitely many of them.

\emph{Back and forth.}  Let \(\mathcal M,\mathcal N\entails T_P\) be
\(\omega\)-saturated, \(f:\bar a\mapsto\bar b\) a finite partial isomorphism,
and \(c\in M\) not in the domain.  The set
\(\Pi(x)=\{P_n^{\sigma_c(n)}(x):n<\omega\}\cup\{x\neq b_i:i\}\) has
parameters from the finite tuple \(\bar b\).  A finite subset of \(\Pi\)
asks for an element with a finite pattern \(P^\sigma\) and different from the
finitely many \(b_i\); the axioms of \(T_P\) give more elements with
\(P^\sigma\) than there are \(b_i\).  So \(\Pi\) is realized in
\(\mathcal N\) by some \(d\).  The atomic formulas are \(P_n(u)\) and
\(u=v\); \(c\) and \(d\) have the same pattern, \(c\neq a_i\) and
\(d\neq b_i\) for all \(i\), so \(f\cup\{(c,d)\}\) is a finite partial
isomorphism.  The other direction is symmetric.  By
\cref{thm:p03-w03-p03-032-qf-test}, \(T_P\) is complete, and for an element
\(c\) and a set \(A\) the type \(\operatorname{tp}(c/A)\) is decided by the
atomic formulas satisfied by \(c\) together with tuples from \(A\), namely by
\(c\) itself if \(c\in A\), and otherwise by the fact \(c\notin A\) and by
\(\sigma_c\).

(2) By (1), \(\abs{S_1(A)}\leq\abs A+2^{\aleph_0}\).  Conversely, distinct
elements of \(A\) have distinct types, and for each \(\sigma\in2^\omega\) the
set \(\{P_n^{\sigma(n)}(x):n<\omega\}\cup\{x\neq a:a\in A\}\) is finitely
satisfiable by the axioms and has small parameter set, so it is realized by
saturation; distinct \(\sigma\) give distinct types.  Hence
\(\abs{S_1(A)}=\abs A+2^{\aleph_0}\).  If \(\lambda\geq2^{\aleph_0}\) and
\(\abs A\leq\lambda\), this is at most \(\lambda\); if
\(\lambda<2^{\aleph_0}\), the set \(A=\varnothing\) has
\(\abs{S_1(\varnothing)}=2^{\aleph_0}>\lambda\).  In particular \(T_P\) is
not \(\aleph_0\)-stable.  It is \(2^{\aleph_0}\)-stable, hence stable by
\cref{thm:p03-w03-p03-032-spectrum}, and since \(\abs{T_P}=\aleph_0\) it is
\(\lambda\)-stable for every \(\lambda\geq2^{\abs{T_P}}\).

(3) \emph{Copies outside a given set.}  Let \(a=(a_1,\ldots,a_r)\) and
\(C\subseteq B\) be as stated, and let \(D\supseteq B\) be small.  Choose
\(a'\) as follows: \(a'_j=a_j\) if \(a_j\in C\); for the entries
\(a_j\notin B\), choose new elements by realizing, by saturation, the set of
formulas in variables \(z_j\) saying that \(z_j\) has pattern
\(\sigma_{a_j}\), that \(z_j=z_l\) iff \(a_j=a_l\), and that
\(z_j\neq e\) for every \(e\in D\); each finite subset is satisfiable by the
axioms of \(T_P\).  Then \((a',\bar e)\) and \((a,\bar e)\) satisfy the same
atomic formulas for every tuple \(\bar e\) from \(B\): patterns agree, entries
in \(C\) are unchanged, and an entry \(a_j\notin B\) differs from every
element of \(B\), as does \(a'_j\).  By \cref{thm:p03-w03-p03-032-qf-test}
\(a'\equiv_B a\), and every entry of \(a'\) that lies in \(D\) lies in \(C\).

\emph{No dividing.}  Let \(\theta(x;b)\in\operatorname{tp}(a/B)\), and let
\((b_i)_{i<\omega}\) be \(C\)-indiscernible with \(b_0=b\).  Apply the
construction with \(D=B\cup\bigcup_i\operatorname{ran}b_i\) to obtain
\(a'\).  For each \(i\), the tuples \((a',b_i)\) and \((a,b)\) satisfy the
same atomic formulas: \(b_i\equiv_C b\), so \(b_i\) and \(b\) satisfy the
same atomic formulas and the same equations \(y_\alpha=e\) with \(e\in C\);
an entry of \(a'\) in \(C\) equals \(b_i(\alpha)\) iff it equals
\(b(\alpha)\); a new entry of \(a'\) lies outside \(D\), so it equals no
\(b_i(\alpha)\), just as the corresponding entry of \(a\), lying outside
\(B\), equals no \(b(\alpha)\); and the patterns agree.  By
\cref{thm:p03-w03-p03-032-qf-test}, \(a'b_i\equiv ab\), so
\(\theta(a';b_i)\) holds for every \(i\), and
\(\{\theta(x;b_i):i<\omega\}\) is consistent.  Hence no formula of
\(\operatorname{tp}(a/B)\) divides over \(C\).  This holds for every
\(B\supseteq C\) and every \(a\) whose entries in \(B\) lie in \(C\).

\emph{No forking.}  Suppose \(\psi(x;e)\in\operatorname{tp}(a/B)\) forks over
\(C\): by \cref{thm:p03-w03-p03-031-saturation}\,(3) there are formulas
\(\varphi_j(x;e_j)\), \(1\leq j\leq k\), each dividing over \(C\), with
\(\mathfrak C\entails\forall x\,(\psi(x;e)\to\bigvee_j\varphi_j(x;e_j))\).
Let \(D=B\cup\operatorname{ran}e_1\cup\cdots\cup\operatorname{ran}e_k\) and
take \(a'\) as above.  Since \(a'\equiv_B a\), \(\psi(a';e)\) holds, so some
\(\varphi_j(a';e_j)\) holds.  The entries of \(a'\) lying in \(D\) lie in
\(C\), so by the previous paragraph applied to \(a'\) and \(D\), no formula of
\(\operatorname{tp}(a'/D)\) divides over \(C\); but
\(\varphi_j(x;e_j)\in\operatorname{tp}(a'/D)\) does.  This contradiction
shows that \(\operatorname{tp}(a/B)\) does not fork over \(C\).  Taking \(C\)
to be the finite set of entries of \(a\) lying in \(B\), and using that
\(T_P\) is stable by (2), \(T_P\) is superstable.
\end{proof}
\end{annexitem}

\begin{annexitem}{annex:p03-w03-p03-032-koenig}{thm:p03-w03-p03-032-koenig}{Countable cofinality raises the power}
\begin{proof}
Let \((g_\alpha)_{\alpha<\lambda}\) be any family of functions
\(\omega\to\lambda\).  For each \(n\), the set
\(\{g_\alpha(n):\alpha<\lambda_n\}\) has at most \(\lambda_n<\lambda\)
members, so there is a least ordinal \(g(n)<\lambda\) outside it.  For each
\(\alpha<\lambda\) some \(n\) has \(\alpha<\lambda_n\), because
\(\lambda=\sup_n\lambda_n\); then \(g(n)\neq g_\alpha(n)\).  So
\(g\neq g_\alpha\) for every \(\alpha\), no family of \(\lambda\) functions
exhausts \(\lambda^\omega\), and \(\lambda^{\aleph_0}>\lambda\).
\end{proof}
\end{annexitem}

\begin{annexitem}{annex:p03-w03-p03-032-te}{thm:p03-w03-p03-032-te}{\(T_E\) is stable and not superstable}
\begin{proof}
(1) \emph{The model \(\mathcal M_\lambda\).}  Each \(E_n\) is an equivalence
relation, and agreement on the first \(n+2\) places implies agreement on
the first \(n+1\).  The \(E_0\)-classes are indexed by the first entry, so
there are \(\lambda\) of them; the \(E_{n+1}\)-classes inside one
\(E_n\)-class are indexed by the entry at place \(n+1\), so there are
\(\lambda\) of them.  So \(\mathcal M_\lambda\entails T_E\).

\emph{Ultrametric law.}  If \(n<\min(h(x,y),h(y,z))\), then
\(x\mathrel{E_n}y\mathrel{E_n}z\), so \(x\mathrel{E_n}z\); hence
\(h(x,z)\geq\min(h(x,y),h(y,z))\).  If \(h(x,y)<h(y,z)\) and
\(h(x,z)>h(x,y)\), then \(h(x,y)\geq\min(h(x,z),h(z,y))>h(x,y)\), which is
impossible; so \(h(x,z)=h(x,y)\).  The case \(h(y,z)<h(x,y)\) is symmetric.

\emph{Atomic formulas.}  These are \(u=v\) and \(u\mathrel{E_n}v\), and
\(u\mathrel{E_n}v\) holds iff \(n<h(u,v)\).  So the atomic formulas satisfied
by a tuple are fixed by its equalities and the heights between its entries,
and finite partial isomorphisms preserve \(h\).

\emph{Back and forth.}  Let \(\mathcal M,\mathcal N\entails T_E\) be
\(\omega\)-saturated, \(f:(a_1,\ldots,a_r)\mapsto(b_1,\ldots,b_r)\) a finite
partial isomorphism and \(c\in M\) outside its domain.  If \(r=0\), any
\(d\in N\) works.  Otherwise let \(K=\max_ih(c,a_i)\), attained at
\(i=j\).  For every \(i\): if \(h(a_j,a_i)<K\), then
\(h(c,a_i)=h(a_j,a_i)\) by the ultrametric law; if \(h(a_j,a_i)\geq K\), then
\(h(c,a_i)\geq\min(K,h(a_j,a_i))=K\), hence \(h(c,a_i)=K\).  So the heights
from \(c\) are determined by \(j\), \(K\) and the heights among the \(a_i\).
It suffices to find \(d\in N\setminus\{b_1,\ldots,b_r\}\) with
\(h(d,b_j)=K\) and \(h(d,b_i)\leq K\) for all \(i\): the same computation
then gives \(h(d,b_i)=h(c,a_i)\) for every \(i\), because
\(h(b_j,b_i)=h(a_j,a_i)\), and \(f\cup\{(c,d)\}\) is a finite partial
isomorphism.
\begin{itemize}
  \item \(K=0\): the \(b_i\) lie in at most \(r\) of the infinitely many
        \(E_0\)-classes; take \(d\) in another one.
  \item \(0<K<\omega\): the \(E_{K-1}\)-class of \(b_j\) contains infinitely
        many \(E_K\)-classes; take \(d\) in one containing no \(b_i\).  Then
        \(d\mathrel{E_{K-1}}b_j\) and \(\neg\,d\mathrel{E_K}b_i\) for all
        \(i\), so \(h(d,b_j)=K\) and \(h(d,b_i)\leq K\).
  \item \(K=\omega\): the set \(\{x\mathrel{E_n}b_j:n<\omega\}\cup
        \{x\neq b_i:i\}\) is finitely satisfiable, since the
        \(E_n\)-class of \(b_j\) contains infinitely many \(E_{n+1}\)-classes
        and so infinitely many elements; by \(\omega\)-saturation some \(d\)
        realizes it, and \(h(d,b_j)=\omega\).
\end{itemize}
The other direction is symmetric.  By \cref{thm:p03-w03-p03-032-qf-test},
\(T_E\) is complete, and the type of a tuple over a set is decided as stated.

(2) \emph{Upper bound.}  Let \(\abs A\leq\lambda\) and \(c\notin A\); by (1)
\(\operatorname{tp}(c/A)\) is decided by the heights \(h(c,a)\),
\(a\in A\).  If \(A\neq\varnothing\), let \(K=\sup_{a\in A}h(c,a)\).  If the
supremum is attained at some \(a^*\), the computation in (1) gives
\(h(c,a)=K\) when \(h(a^*,a)\geq K\) and \(h(c,a)=h(a^*,a)\) otherwise, so
the type is decided by the pair \((a^*,K)\in A\times(\omega+1)\).  If it is
not attained, then \(K=\omega\) and every \(h(c,a)\) is finite; choose
\(a_n\in A\) with \(h(c,a_n)>n\).  For \(a\in A\) take \(n=h(c,a)\): then
\(h(c,a_n)>h(c,a)\), so \(h(c,a)=h(a_n,a)\) by the ultrametric law, and the
type is decided by \((a_n)_{n<\omega}\in A^\omega\).  Counting also the
types of members of \(A\) and the single type over \(\varnothing\),
\(\abs{S_1(A)}\leq\abs A+\abs A\cdot\aleph_0+\abs A^{\aleph_0}+1
\leq\lambda^{\aleph_0}\).

\emph{Lower bound.}  In \(\mathcal M_\lambda\) let
\(A_0=\{s^\frown0^\omega:s\in\lambda^{<\omega}\}\), the sequences that are
eventually \(0\); \(\abs{A_0}=\lambda\).  If \(\eta\neq\eta'\) in
\(\lambda^\omega\) first differ at place \(n\), let
\(e=(\eta\mathbin{\upharpoonright}(n+1))^\frown0^\omega\in A_0\); then
\(\eta\mathrel{E_n}e\) and \(\neg\,\eta'\mathrel{E_n}e\).  So the elements of
\(\lambda^\omega\) have \(\lambda^{\aleph_0}\) distinct types over \(A_0\) in
\(\mathcal M_\lambda\).  Let \(\Delta\) be the set of formulas
\(\theta(z_{e_1},\ldots,z_{e_k})\), in variables indexed by \(A_0\), with
\(\mathcal M_\lambda\entails\theta(e_1,\ldots,e_k)\).  Each finite subset of
\(\Delta\) is implied by one existential sentence true in
\(\mathcal M_\lambda\), hence true in \(\mathfrak C\) because \(T_E\) is
complete; so \(\Delta\) is realized by saturation, by a set
\(A=\{e':e\in A_0\}\) with \((\mathfrak C,e')_{e\in A_0}\equiv
(\mathcal M_\lambda,e)_{e\in A_0}\).  Each of the \(\lambda^{\aleph_0}\)
distinct complete types over \(A_0\) realized in \(\mathcal M_\lambda\)
becomes, after renaming \(e\) as \(e'\), a set of formulas over \(A\) each of
whose finite parts is satisfied in \(\mathfrak C\); realized by saturation,
these give \(\lambda^{\aleph_0}\) distinct members of \(S_1(A)\).

\emph{Consequences.}  If \(\lambda^{\aleph_0}=\lambda\), then
\(T_E\) is \(\lambda\)-stable by the upper bound; if
\(\lambda^{\aleph_0}>\lambda\), the lower bound shows it is not.  As
\((2^{\aleph_0})^{\aleph_0}=2^{\aleph_0}\), \(T_E\) is
\(2^{\aleph_0}\)-stable and so stable by
\cref{thm:p03-w03-p03-032-spectrum}.  The cardinal
\(\beth_\omega=\sup_n\beth_n\) has \(\beth_n<\beth_\omega\), so
\(\beth_\omega^{\aleph_0}>\beth_\omega\) by
\cref{thm:p03-w03-p03-032-koenig}, and \(T_E\) is not
\(\beth_\omega\)-stable, while \(\beth_\omega\geq\beth_1=2^{\aleph_0}
=2^{\abs{T_E}}\).

(3) \emph{The configuration.}  In \(\mathcal M_{\aleph_0}\) let \(a_n\) be
the sequence with \(1\) at the places below \(n\) and \(0\) elsewhere, and
\(c\) the constant sequence \(1\).  For \(n<m\), \(a_n\) and \(a_m\) first
differ at place \(n\), and \(c\) and \(a_n\) first differ at place \(n\).  So
\(h(a_n,a_m)=\min(n,m)\) and \(h(c,a_n)=n\).  As in the lower bound, the
complete type of \((c,a_0,a_1,\ldots)\) in \(\mathcal M_{\aleph_0}\) is
finitely satisfiable in \(\mathfrak C\) and is realized there; we use the
realization and keep the names.

\emph{Forking.}  Let \(B_0\subseteq B\) be finite, and choose \(N\) with
\(B_0\subseteq\{a_0,\ldots,a_N\}\).  Put \(n=N+2\) and
\(\theta(x;a_n)=x\mathrel{E_{n-1}}a_n\); since \(n-1<n=h(c,a_n)\),
\(\theta(x;a_n)\in\operatorname{tp}(c/B)\).  The \(E_{n-2}\)-class of
\(a_n\) contains infinitely many \(E_{n-1}\)-classes; let \(e_0=a_n\) and
choose \(e_i\) for \(0<i<\omega\) in that \(E_{n-2}\)-class, in pairwise
distinct \(E_{n-1}\)-classes all different from that of \(a_n\).  Then
\(h(e_i,e_k)=n-1\) for \(i\neq k\).  For \(a_l\in B_0\) we have \(l\leq N<n-1\),
so \(h(a_n,a_l)=l<n-1=h(e_i,a_n)\) for \(i>0\), and the ultrametric law gives
\(h(e_i,a_l)=l=h(a_n,a_l)\).  Hence any two increasing \(r\)-tuples from
\((e_i)_{i<\omega}\) satisfy the same atomic formulas together with any tuple
from \(B_0\): all entries are distinct and outside \(B_0\) (their heights to
members of \(B_0\) are finite), all heights among them are \(n-1\), and the
heights to \(a_l\in B_0\) are \(l\).  By
\cref{thm:p03-w03-p03-032-qf-test} the sequence is \(B_0\)-indiscernible.
The set \(\{\theta(x;e_0),\theta(x;e_1)\}\) is inconsistent, since
\(x\mathrel{E_{n-1}}e_0\) and \(x\mathrel{E_{n-1}}e_1\) would give
\(e_0\mathrel{E_{n-1}}e_1\).  So \(\theta(x;a_n)\) divides over \(B_0\), and
therefore forks over \(B_0\) (take \(m=1\) and \(\varphi_1=\theta\) in
\cref{def:p03-dividing-forking}).  Thus \(\operatorname{tp}(c/B)\) forks over
every finite subset of \(B\), and the stable theory \(T_E\) is not
superstable.
\end{proof}
\end{annexitem}
